
Conversational AI systems are trained to be helpful through alignment methods such as Reinforcement Learning from Human Feedback (RLHF) (Ouyang et al., 2022; Bai et al., 2022), which optimize AI policy until human evaluators rate outputs favorably. This approach treats alignment as unidirectional optimization: adjust the AI until humans are satisfied. It does not account for the bidirectional nature of conversation observed in human-human dialogue, where both participants adapt to each other over time \citep{pickering2004toward}.

During multi-turn conversations with AI assistants, users are not static evaluators. They refine their communication strategies based on what works: reformulating unsuccessful queries, learning which requests the AI handles well, and developing mental models of AI capabilities. Simultaneously, RLHF-trained systems may exhibit responsiveness to user query patterns through learned behaviors. This bidirectional co-adaptation suggests that alignment is not solely a property of AI policy but an emergent property of the conversation itself.

Current alignment evaluation methods do not capture this interactive dynamic. RLHF optimizes for human preference accuracy, evaluating single-turn quality or aggregate satisfaction scores rather than how users and AI co-adapt within conversations. Metrics that distinguish conversations where genuine bidirectional alignment occurs---users learn to communicate effectively AND the AI exhibits responsiveness---are absent.

This paper proposes behavioral coupling as a metric for bidirectional alignment. The key insight is that co-adaptation produces observable behavioral signatures: users who learn AI capabilities reformulate queries less over turns (learning curve effect), while responsive AI systems produce outputs whose diversity correlates with query diversity. By measuring correlation between these two signals---user learning rate (reformulation slope) and AI responsiveness (query-response diversity correlation)---we quantify bidirectional alignment strength.

To validate this framework, 169,352 multi-turn conversations from the HH-RLHF dataset \citep{bai2022constitutional} were analyzed, focusing on conversations with $\geq 5$ turns where learning curves can be detected. Two existence hypotheses were tested: (1) users exhibit decreasing reformulation rates over turns, and (2) AI response diversity correlates with query diversity. Results show:

\begin{itemize}
\item User learning exists: Reformulation rates decrease over conversation turns (mean slope = -0.021, p=0.012, Cohen's d=-0.246).
\item AI responsiveness exists: Response diversity correlates with query diversity (r=0.396, p<0.001, 95\% CI [0.392, 0.400]).
\item Temporal co-adaptation: Diversity correlation strengthens with conversation length (r=0.04 for 2-3 turns to r=0.19 for 6+ turns).
\end{itemize}

Effect sizes are smaller than initially hypothesized (small Cohen's d for learning, correlation below the preregistered r>0.4 threshold), but both components are statistically robust across a large sample. These findings validate the theoretical framework's foundation. However, the core hypothesis---whether coupling strength predicts conversation quality beyond individual metrics---remains untested.

This work makes three contributions: (1) It introduces behavioral coupling as an operationalization of bidirectional alignment, shifting measurement from static policy evaluation to conversation dynamics. (2) It provides empirical evidence that both components (user learning and AI responsiveness) exist in RLHF-trained systems. (3) It demonstrates that these behavioral signals can be extracted from conversation logs without additional human evaluation, enabling scalable alignment monitoring.

The remainder of the paper is organized as follows: Section 2 reviews related work. Section 3 describes methodology for detecting reformulation patterns and measuring diversity correlations. Section 4 presents experimental design and datasets. Section 5 reports results. Section 6 discusses implications and limitations. Section 7 concludes.
